% !TEX program = xelatex
% 梯度下降：从零到收敛的完整推导
% 编译：xelatex gradient_descent.tex（跑两遍以生成目录与书签）
\documentclass[UTF8, a4paper, zihao=-4, fontset=none]{ctexart}

% ---------------------------------------------------------------- 字体
% Linux 上 Noto CJK 是 ttc 合集，按字体名匹配会落到日文 face（JP），
% 必须用 FontIndex 指定简体中文 face；其它系统退回字体名匹配。
\IfFileExists{/usr/share/fonts/opentype/noto/NotoSerifCJK-Regular.ttc}{%
  \setCJKmainfont{NotoSerifCJK-Regular.ttc}[Path=/usr/share/fonts/opentype/noto/, FontIndex=2,
    BoldFont=NotoSerifCJK-Bold.ttc, BoldFeatures={FontIndex=2}]%
  \setCJKsansfont{NotoSansCJK-Regular.ttc}[Path=/usr/share/fonts/opentype/noto/, FontIndex=2,
    BoldFont=NotoSansCJK-Bold.ttc, BoldFeatures={FontIndex=2}]%
  \setCJKmonofont{NotoSansCJK-Regular.ttc}[Path=/usr/share/fonts/opentype/noto/, FontIndex=7,
    BoldFont=NotoSansCJK-Bold.ttc, BoldFeatures={FontIndex=7}]%
}{%
  \IfFontExistsTF{Noto Serif CJK SC}{%
    \setCJKmainfont{Noto Serif CJK SC}%
    \setCJKsansfont{Noto Sans CJK SC}%
    \setCJKmonofont{Noto Sans Mono CJK SC}%
  }{%
    % Windows/macOS：优先系统 TrueType 中文字体。Fandol / Noto CJK 属于
    % CID-keyed CFF，PDF 会依赖 Adobe-GB1 CMap，部分阅读器（含 poppler 预览）不显示中文。
    \IfFontExistsTF{SimSun}{%
      \setCJKmainfont{SimSun}[BoldFont=simhei.ttf, AutoFakeSlant=0.2]%
      \setCJKsansfont{Microsoft YaHei}[BoldFont={Microsoft YaHei Bold}]%
      \setCJKmonofont{SimSun}[AutoFakeSlant=0.2]%
    }{%
      \setCJKmainfont{FandolSong-Regular.otf}[BoldFont=FandolSong-Bold.otf]%
      \setCJKsansfont{FandolHei-Regular.otf}[BoldFont=FandolHei-Bold.otf]%
      \setCJKmonofont{FandolSong-Regular.otf}%
    }%
  }%
}
% 圈数字（①②…）交给中文字体渲染，否则拉丁字体缺字
\xeCJKDeclareCharClass{CJK}{"2460 -> "24FF}
\xeCJKDeclareCharClass{CJK}{"25A0 -> "25FF}

% ---------------------------------------------------------------- 宏包
\usepackage[margin=2.5cm, top=2.6cm, bottom=2.6cm]{geometry}
\usepackage{amsmath, amssymb, mathtools, bm}
\usepackage{booktabs, tabularx, array, multirow}
\usepackage{caption}
\usepackage[most]{tcolorbox}
\usepackage{listings}
\usepackage{xcolor}
\usepackage{enumitem}
\usepackage{fancyhdr}
\usepackage{titlesec}
\usepackage{lastpage}
\usepackage{needspace}
\usepackage[hidelinks, bookmarksnumbered]{hyperref}

% ---------------------------------------------------------------- 配色
\definecolor{accent}{HTML}{1F4E79}
\definecolor{accentsoft}{HTML}{EAF1F7}
\definecolor{warn}{HTML}{9C3A2A}
\definecolor{warnsoft}{HTML}{FBEEEA}
\definecolor{good}{HTML}{1D6B4F}
\definecolor{goodsoft}{HTML}{E9F4EF}
\definecolor{rulegray}{HTML}{E2DDD4}
\definecolor{codebg}{HTML}{F6F4EF}

% ---------------------------------------------------------------- 版式
\linespread{1.32}
\setlength{\emergencystretch}{2em}
\setlength{\parskip}{0.35em}
\setlist[itemize]{leftmargin=1.6em, itemsep=0.15em, topsep=0.3em}
\setlist[enumerate]{leftmargin=1.8em, itemsep=0.15em, topsep=0.3em}
\captionsetup{font=small, labelfont=bf, textfont=small, skip=4pt}
\setlength{\parindent}{2em}

\titleformat{\section}
  {\normalfont\Large\bfseries\color{accent}}
  {\thesection}{0.6em}{}[\vspace{-0.55em}\color{rulegray}\rule{\linewidth}{0.6pt}]
\titlespacing*{\section}{0pt}{2.1em}{0.9em}
\titleformat{\subsection}
  {\normalfont\large\bfseries}
  {\thesubsection}{0.5em}{}
\titlespacing*{\subsection}{0pt}{1.4em}{0.5em}

\setcounter{tocdepth}{1}
\renewcommand{\contentsname}{目录}

\pagestyle{fancy}
\fancyhf{}
\fancyfoot[L]{\footnotesize\color{gray}梯度下降：从零到收敛的完整推导}
\fancyfoot[R]{\footnotesize\color{gray}\thepage\,/\,\pageref{LastPage}}
\renewcommand{\headrulewidth}{0pt}

% ---------------------------------------------------------------- 盒子
\newtcolorbox{infobox}{enhanced, breakable, pad at break*=4pt,
  colback=accentsoft, colframe=accentsoft, boxrule=0pt,
  borderline west={3pt}{0pt}{accent}, arc=1pt, left=10pt, right=10pt,
  top=8pt, bottom=8pt}
\newtcolorbox{warnbox}{enhanced, breakable, pad at break*=4pt,
  colback=warnsoft, colframe=warnsoft, boxrule=0pt,
  borderline west={3pt}{0pt}{warn}, arc=1pt, left=10pt, right=10pt,
  top=8pt, bottom=8pt}
\newtcolorbox{goodbox}{enhanced, breakable, pad at break*=4pt,
  colback=goodsoft, colframe=goodsoft, boxrule=0pt,
  borderline west={3pt}{0pt}{good}, arc=1pt, left=10pt, right=10pt,
  top=8pt, bottom=8pt}
\newtcolorbox{examplebox}{enhanced, breakable, pad at break*=4pt,
  colback=white, colframe=rulegray, boxrule=0.5pt, arc=1pt,
  left=10pt, right=10pt, top=7pt, bottom=7pt}
\newtcolorbox{answerbox}{enhanced, breakable, pad at break*=4pt,
  colback=white, colframe=rulegray, boxrule=0.5pt, arc=1pt,
  left=10pt, right=10pt, top=7pt, bottom=7pt}

% ---------------------------------------------------------------- 代码
\lstset{
  language=Python,
  basicstyle=\ttfamily\footnotesize,
  keywordstyle=\color{accent}\bfseries,
  commentstyle=\color{gray!80!black}\itshape,
  stringstyle=\color{good},
  backgroundcolor=\color{codebg},
  showstringspaces=false,
  breaklines=true,
  columns=fullflexible,
  keepspaces=true,
  frame=none,
  xleftmargin=6pt, xrightmargin=6pt,
  aboveskip=8pt, belowskip=8pt,
  extendedchars=true
}

\hypersetup{
  pdftitle={梯度下降：从零到收敛的完整推导},
  pdfsubject={机器学习基础讲义},
  pdfcreator={XeLaTeX + ctex}
}

\begin{document}


\begin{center}
{\sffamily\small\color{accent}\bfseries 机器学习基础 · 讲义}\\[6pt]
{\Huge\bfseries 梯度下降：从零到收敛的完整推导}\\[4pt]
{\color{rulegray}\rule{0.55\linewidth}{0.6pt}}
\end{center}
\vspace{1.2em}

\tableofcontents
\newpage





\needspace{8\baselineskip}
\section{把“学习”翻译成“最小化”}




\needspace{5\baselineskip}
\subsection{三个东西：参数、损失、目标}



训练一个模型，说到底只做一件事：\textbf{找一组参数，让某个衡量“错得多离谱”的数尽可能小}。这句话里其实有三个对象，先把它们分清楚，后面的推导才有落脚点。



\begin{center}
\footnotesize
\setlength{\tabcolsep}{3pt}
\begin{tabularx}{\linewidth}{@{}>{\hsize=0.545\hsize\raggedright\arraybackslash}X>{\hsize=0.545\hsize\raggedright\arraybackslash}X>{\hsize=1.778\hsize\raggedright\arraybackslash}X>{\hsize=1.132\hsize\raggedright\arraybackslash}X@{}}
\toprule
\textbf{对象} & \textbf{记号} & \textbf{含义} & \textbf{在直线拟合里是什么} \\
\midrule
参数 & $\mathbf{w}\in\mathbb{R}^d$ & 模型里可以被调整的数，$d$ 是它的个数 & 斜率，若带截距则还有 $b$ \\
损失函数 & $L(\mathbf{w})$ & 给定一组参数时，模型错得有多严重；越小越好 & 所有点的预测偏差平方和 \\
目标 & $\mathbf{w}^\star=\arg\min_{\mathbf{w}} L(\mathbf{w})$ & 让损失取最小的那组参数 & 最小二乘解 \\
\bottomrule
\end{tabularx}
\end{center}



注意 $L$ 是一个\textbf{从 $\mathbb{R}^d$ 到 $\mathbb{R}$ 的函数}：你喂给它 $d$ 个数，它还你一个数。梯度下降要做的，是在这个 $d$ 维的“地形”上，从一个初始点出发，一步步走到谷底。



\begin{infobox}
{\bfseries\sffamily\small\color{accent} 先建立一个图像}\par\vspace{2pt}
把 $L(\mathbf{w})$ 想成一片山地，$\mathbf{w}$ 是你站在地图上的坐标（经度、纬度……共 $d$ 个方向），$L$ 是你脚下的海拔。你的目标是走到最低的谷底。但你是个\textbf{蒙着眼睛的人}：看不见全局，只能用脚感受脚下这一小块的坡度。梯度下降就是“摸一摸脚下往哪边最陡，就往反方向迈一步”的规则。


这个比喻会把后面所有结论都解释清楚：为什么只能保证局部、为什么步长太大会摔下悬崖、为什么地形狭长时会来回横跳。


\end{infobox}




\needspace{5\baselineskip}
\subsection{一个能算到底的例子}



抽象的话说到这里，先看一个能手算的例子，它在第 7 节会被完整走一遍。



假设有三个观测点 $(x,y)=(1,\,2.1),\ (2,\,3.9),\ (3,\,6.2)$，我们想用一条过原点的直线 $\hat{y}=wx$ 去拟合。参数只有一个：斜率 $w$。



预测偏差是 $\hat{y}_i-y_i=wx_i-y_i$，把三个偏差的平方加起来（平方是为了让正负偏差都算“错”，且对大偏差惩罚更重），再除以 2 只是为了后面求导时消掉系数：


\[
L(w)=\frac{1}{2}\sum_{i=1}^{3}(wx_i-y_i)^2
=\frac12\Big[(w-2.1)^2+(2w-3.9)^2+(3w-6.2)^2\Big]
\]


展开合并同类项（你可以自己验一遍）：


\[
L(w)=7w^2-28.5w+29.03
\]


这是一条开口向上的抛物线。它的最小值在哪？对 $w$ 求导令其为零：


\[
L'(w)=14w-28.5=0\quad\Longrightarrow\quad w^\star=\frac{28.5}{14}\approx 2.036
\]


答案精确可得，因为这是一维二次函数。但真实模型里 $L$ 可能是几亿个参数的复杂函数，\textbf{没有公式解}，只能一步步逼近。梯度下降就是那个“一步步逼近”的算法。



\begin{goodbox}
{\bfseries\sffamily\small\color{good} 这一节要记住的}\par\vspace{2pt}
梯度下降不是“另一种求最小值的方法”，而是在\textbf{没有解析解}时唯一可行的通用手段。它的全部信息来自“当前点的局部斜率”，所以它的所有优点和缺陷，都源于这个信息限制。


\end{goodbox}




\needspace{8\baselineskip}
\section{一维：导数的符号告诉你往哪走}




\needspace{5\baselineskip}
\subsection{导数究竟说了什么}



一元函数 $f(w)$ 在 $w_0$ 处的导数定义为


\[
f'(w_0)=\lim_{h\to 0}\frac{f(w_0+h)-f(w_0)}{h}
\]


这个式子的意思是：\textbf{在 $w_0$ 处把 $w$ 增大一点点，函数值会以多大的速率变化}。它是一个“变化率”，不是一个“位置”。



由此直接得到两条判据，这是整篇文章的地基：



\begin{itemize}
\item $f'(w_0)>0$：往右走，函数\textbf{上升}。所以想下降，就该往\textbf{左}走，即减小 $w$。
\item $f'(w_0)<0$：往右走，函数\textbf{下降}。所以想下降，就该往\textbf{右}走，即增大 $w$。
\end{itemize}



两种情况能不能写成一个式子？可以。取


\[
w_1=w_0-\eta f'(w_0),\qquad \eta>0
\]


当 $f'(w_0)>0$ 时，减去一个正数，$w$ 变小，往左走，正确；当 $f'(w_0)<0$ 时，减去一个负数等于加上一个正数，$w$ 变大，往右走，也正确。\textbf{负号一次性处理了两种情形}，这就是“减”的来历，它不是什么深奥的规定，只是把“顺着坡往下”这件事写成了一个统一公式。




\needspace{5\baselineskip}
\subsection{为什么“减一点”就一定会下降}



上面的论证还是定性的。要定量知道“走多远、降多少”，需要泰勒展开。



若 $f$ 在 $w_0$ 处可导，则


\[
f(w_0+h)=f(w_0)+f'(w_0)\,h+o(h)
\]


其中 $o(h)$ 表示“比 $h$ 更快趋于零的项”，即 $\lim_{h\to0}o(h)/h=0$。这句话的含义是：\textbf{当 $h$ 足够小时，函数的变化几乎完全由线性项 $f'(w_0)h$ 决定}，剩下的部分是更高阶的小量。



代入 $h=-\eta f'(w_0)$：


\[
f\big(w_0-\eta f'(w_0)\big)=f(w_0)-\eta\big[f'(w_0)\big]^2+o(\eta)
\]


关键在 $-\eta[f'(w_0)]^2$ 这一项：只要 $f'(w_0)\neq0$ 且 $\eta>0$，它就\textbf{严格为负}。而 $o(\eta)$ 在 $\eta$ 足够小时可以任意小，小到无法抵消这一项。于是存在某个 $\delta>0$，使得所有 $0<\eta<\delta$ 都满足


\[
f\big(w_0-\eta f'(w_0)\big)<f(w_0)
\]


\begin{infobox}
{\bfseries\sffamily\small\color{accent} 严格性的来源}\par\vspace{2pt}
上面这一步不是“近似”，而是严格结论。它的逻辑是：\emph{极限为正 \ensuremath{\Longrightarrow} 存在邻域使函数值为正}。把导数定义写成


取 $h=-\eta f'(w_0)$，则当 $\eta\to0^+$ 时，比值的极限是 $-f'(w_0)^2<0$。按极限定义，存在 $\delta$，使 $0<\eta<\delta$ 时该比值为负，即分子 $f(w_0+h)-f(w_0)<0$。这就证明了“小步必然下降”。


注意结论里的“存在 $\delta$”：\textbf{下降只在足够小的步长下成立}，这个限定会在第 4 节变成具体的数值上界。


\end{infobox}




\needspace{5\baselineskip}
\subsection{步长的三种命运}



$\eta$ 叫学习率（learning rate）或步长。它决定了沿着斜率方向迈多远。用同一个抛物线 $L(w)=7w^2-28.5w+29.03$（$L''=14$）从 $w_0=0$ 出发，可以清楚看到三种结局：



\begin{center}
\footnotesize
\setlength{\tabcolsep}{3pt}
\begin{tabularx}{\linewidth}{@{}>{\hsize=0.560\hsize\raggedleft\arraybackslash}X>{\hsize=0.925\hsize\raggedleft\arraybackslash}X>{\hsize=0.888\hsize\raggedright\arraybackslash}X>{\hsize=1.627\hsize\raggedright\arraybackslash}X@{}}
\toprule
\textbf{$\eta$} & \textbf{前几步的 $w$} & \textbf{现象} & \textbf{原因} \\
\midrule
0.02 & 0.57 \ensuremath{\rightarrow} 0.98 \ensuremath{\rightarrow} 1.28 \ensuremath{\rightarrow} 1.49 \ensuremath{\rightarrow} … & 稳稳逼近 2.036，但很慢 & 每步只走一点点，误差按固定比例缩小 \\
0.12 & 3.42 \ensuremath{\rightarrow} 1.09 \ensuremath{\rightarrow} 2.68 \ensuremath{\rightarrow} 1.60 \ensuremath{\rightarrow} … & 来回摆动，摆幅逐渐变小 & 一步迈过了谷底，但还没越过对面的坡 \\
0.15 & 4.28 \ensuremath{\rightarrow} −0.43 \ensuremath{\rightarrow} 4.75 \ensuremath{\rightarrow} −0.95 \ensuremath{\rightarrow} … & 越摆越大，发散 & 每步都跨到比出发点更高的位置，能量越积越多 \\
\bottomrule
\end{tabularx}
\end{center}
{\footnotesize\color{gray}\raggedright 初值 $w_0=0$，目标 $w^\star=2.036$。$\eta=0.143$ 是临界点，恰好 $2/L''$。\par}



为什么临界点正好是 $2/14\approx0.143$？这个数字不是巧合。对二次函数，我们可以把每步的误差写成一个递推。



令 $e_k=w_k-w^\star$ 表示“离谷底还有多远”。因为 $L'(w)=14(w-w^\star)=14e_k$，更新式变成


\[
e_{k+1}=w_{k+1}-w^\star=w_k-\eta\cdot14e_k-w^\star=(1-14\eta)\,e_k
\]


于是 $e_k=(1-14\eta)^k e_0$。要让误差趋于零，必须


\[
|1-14\eta|<1\quad\Longleftrightarrow\quad 0<\eta<\frac{2}{14}\approx0.143
\]


\begin{itemize}
\item $\eta<1/14$：$0<1-14\eta<1$，误差单调缩小，\textbf{不摆动}；
\item $1/14<\eta<2/14$：$-1<1-14\eta<0$，误差缩小但每步变号，\textbf{来回摆}；
\item $\eta>2/14$：$|1-14\eta|>1$，误差被放大，\textbf{发散}。
\end{itemize}



这三条正好对应表里的三行。\textbf{“步长必须小于 $2/\text{曲率}$”是梯度下降最基本的一条纪律}，第 4 节会把它推广到多维。



\begin{warnbox}
{\bfseries\sffamily\small\color{warn} 一个容易误解的点}\par\vspace{2pt}
「学习率越小越安全」只对了一半。$\eta$ 太小确实不会发散，但收敛会慢到不可接受；而且在实际训练中，过小的步长会让参数卡在梯度接近零的平坦区域，误以为已经收敛。$\eta$ 的选择是\textbf{稳定性与速度的权衡}，不是单调的“越小越好”。


\end{warnbox}




\needspace{8\baselineskip}
\section{多维：从偏导到梯度}




\needspace{5\baselineskip}
\subsection{偏导：把其他变量按住不动}



模型几乎不会只有一个参数。以 $L(w_1,w_2)$ 为例，它是一个曲面。站在某一点，脚下有无数个方向可以走，每个方向的坡度都可能不同。



先看最简单的一类方向：\textbf{沿着某根坐标轴走}。只让 $w_1$ 变、$w_2$ 不动，此时函数对 $w_1$ 的变化率叫偏导数：


\[
\frac{\partial L}{\partial w_1}=\lim_{h\to0}\frac{L(w_1+h,\ w_2)-L(w_1,w_2)}{h}
\]


它的计算方法和一元导数完全一样：\textbf{把 $w_2$ 当成常数}。例如


\[
L(w_1,w_2)=\frac12\left(w_1^2+3w_2^2\right)
\quad\Longrightarrow\quad
\frac{\partial L}{\partial w_1}=w_1,\qquad
\frac{\partial L}{\partial w_2}=3w_2
\]


偏导告诉你“只动这一个参数时，往哪边、以多快的速率变化”。但它\textbf{不能直接拿来更新}：如果分别用 $-\eta\partial L/\partial w_1$ 和 $-\eta\partial L/\partial w_2$ 去改两个参数，你其实是在同时沿两个坐标轴走，而“同时走”的合成方向是不是最陡的？这需要先回答一个更基本的问题。




\needspace{5\baselineskip}
\subsection{方向导数：任意方向的坡度}



设 $\mathbf{d}$ 是一个单位向量（$\|\mathbf{d}\|=1$），表示一个方向。从点 $\mathbf{w}$ 沿 $\mathbf{d}$ 走，函数的变化率叫\textbf{方向导数}：


\[
D_{\mathbf{d}}L(\mathbf{w})=\lim_{t\to0^+}\frac{L(\mathbf{w}+t\mathbf{d})-L(\mathbf{w})}{t}
\]


由多元链式法则（把 $L(\mathbf{w}+t\mathbf{d})$ 看成 $t$ 的一元函数，对 $t$ 求导）：


\[
\frac{d}{dt}L(\mathbf{w}+t\mathbf{d})\Big|_{t=0}
=\sum_{j=1}^{d}\frac{\partial L}{\partial w_j}\cdot\frac{d(w_j+td_j)}{dt}
=\sum_{j=1}^{d}\frac{\partial L}{\partial w_j}d_j
=\nabla L(\mathbf{w})^\top\mathbf{d}
\]


右边出现的这个向量，就是\textbf{梯度}：把所有偏导按顺序排成一列。


\[
\nabla L(\mathbf{w})=\begin{pmatrix}\dfrac{\partial L}{\partial w_1}\\[6pt]\vdots\\[6pt]\dfrac{\partial L}{\partial w_d}\end{pmatrix}
\]


\begin{goodbox}
{\bfseries\sffamily\small\color{good} 这个恒等式是全部推导的枢纽}\par\vspace{2pt}
\[
D_{\mathbf{d}}L(\mathbf{w})=\nabla L(\mathbf{w})^\top\mathbf{d}
\]


它说的是：\textbf{任意方向的坡度，等于梯度与该方向的内积}。左边是“几何”（沿某个方向走有多陡），右边是“代数”（两个向量的夹角余弦乘以长度）。整个梯度下降的理论，就是反复使用这个等式。


\end{goodbox}




\needspace{5\baselineskip}
\subsection{哪个方向最陡}



现在可以回答“往哪走下降最快”了。要最小化 $D_{\mathbf{d}}L=\nabla L^\top\mathbf{d}$，约束是 $\|\mathbf{d}\|=1$。用 Cauchy–Schwarz 不等式：


\[
\nabla L^\top\mathbf{d}\ \ge\ -\|\nabla L\|\,\|\mathbf{d}\|\ =\ -\|\nabla L\|
\]


等号成立的条件是 $\mathbf{d}$ 与 $\nabla L$ \textbf{反向且共线}，即 $\mathbf{d}=-\nabla L/\|\nabla L\|$。代回去验证：


\[
\nabla L^\top\left(-\frac{\nabla L}{\|\nabla L\|}\right)=-\frac{\|\nabla L\|^2}{\|\nabla L\|}=-\|\nabla L\|
\]


所以：\textbf{负梯度方向是单位步长下函数下降最快的方向}，“最速下降法”（steepest descent）这个名字就是这么来的。



\begin{warnbox}
{\bfseries\sffamily\small\color{warn} "最快"的准确含义}\par\vspace{2pt}
这里的最快是在\textbf{无穷小步长}意义下、并且用\textbf{欧氏长度}衡量“步长”时成立的。换成有限步长，“负梯度方向下降最多”就可能不成立，第 5 节的二次函数例子会给出反例。换成别的长度定义（比如每个坐标单独限幅），最速方向也不再是负梯度。这个限定非常关键，第 10 节会专门讲。


\end{warnbox}




\needspace{5\baselineskip}
\subsection{为什么“负梯度”能同时照顾所有参数}



上面的证明是代数的，但直觉上“为什么偏偏是它”值得单独说。把一阶变化按参数拆开：


\[
\Delta L\approx\sum_{j=1}^{d}\frac{\partial L}{\partial w_j}\Delta w_j
\]


每一项是“第 $j$ 个参数这一步带来的变化”。现在取 $\Delta w_j=-\eta\,\partial L/\partial w_j$，每一项变成


\[
\frac{\partial L}{\partial w_j}\Delta w_j=-\eta\left(\frac{\partial L}{\partial w_j}\right)^2\le 0
\]


\textbf{每一项都非正，没有任何一个参数在往上拉。} 这是负梯度方向独有的性质。换任何别的方向，必然有某些 $\Delta w_j$ 与 $\partial L/\partial w_j$ 同号，那些参数就在“帮倒忙”，需要靠其他参数的下降来补偿，净效果只会更差。



\begin{examplebox}
{\bfseries\sffamily\small\color{black} 数值验证}\par\vspace{2pt}
取 $L=\frac12\|\mathbf{w}\|^2$，点 $\mathbf{w}=(1,3)$，则 $\nabla L=(1,3)$，$\|\nabla L\|=\sqrt{10}\approx3.162$。固定步长 $\|\Delta\mathbf{w}\|=0.5$，把方向从 $-\nabla L$ 开始逐渐旋转，观察两项贡献：


\begin{center}
\footnotesize
\setlength{\tabcolsep}{3pt}
\begin{tabularx}{\linewidth}{@{}>{\hsize=1.250\hsize\raggedright\arraybackslash}X>{\hsize=1.050\hsize\raggedright\arraybackslash}X>{\hsize=0.900\hsize\raggedleft\arraybackslash}X>{\hsize=0.900\hsize\raggedleft\arraybackslash}X>{\hsize=0.900\hsize\raggedleft\arraybackslash}X@{}}
\toprule
\textbf{与 $-\nabla L$ 夹角} & \textbf{$\Delta\mathbf{w}$} & \textbf{$w_1$ 项} & \textbf{$w_2$ 项} & \textbf{净 $\Delta L$} \\
\midrule
0° & $(-0.158,\,-0.474)$ & −0.158 & −1.423 & \textbf{−1.456} \\
45° & $(-0.447,\,-0.224)$ & −0.447 & −0.671 & −0.993 \\
90° & $(-0.474,\,+0.158)$ & −0.474 & +0.474 & +0.125 \\
135° & $(-0.224,\,+0.447)$ & −0.224 & +1.342 & +1.243 \\
180° & $(+0.158,\,+0.474)$ & +0.158 & +1.423 & +1.706 \\
\bottomrule
\end{tabularx}
\end{center}
{\footnotesize\color{gray}\raggedright 第一行两项都是负的，净下降最多；角度一大，$w_2$ 那项翻正并抵消 $w_1$ 的收益；90° 时恰好抵消（一阶为零，只剩二阶项 $+0.125$）。\par}


\end{examplebox}




\needspace{5\baselineskip}
\subsection{梯度与等高线垂直}



等高线是“函数值相同的点的集合”。沿等高线的切线方向 $\mathbf{d}$ 走，函数值不变，所以 $D_{\mathbf{d}}L=\nabla L^\top\mathbf{d}=0$。这说明\textbf{梯度与等高线的切线正交}。



把两件事合起来看：梯度垂直于等高线、指向函数值增大的那一侧，负梯度则垂直指向减小的一侧。于是“梯度下降”在地图上就是\textbf{每一步都垂直穿过等高线，朝更低的那条等高线走}。这也解释了为什么在狭长的山谷里它会左右横跳：垂直穿过等高线的方向，恰好是山谷的横向，而不是沿着谷底延伸的方向。



\begin{infobox}
{\bfseries\sffamily\small\color{accent} 一个小结}\par\vspace{2pt}
到这里，更新公式的每一个零件都有了理由：\newline 
  $\mathbf{w}_{k+1}=\mathbf{w}_k-\eta\nabla L(\mathbf{w}_k)$\newline 
\textbf{减号}：让方向与梯度相反（3.3 的 Cauchy–Schwarz）；\newline 
\textbf{梯度}：唯一让所有参数同向下坡的方向（3.4）；\newline 
\textbf{$\eta$}：控制走多远，必须小到一阶近似仍然有效（2.3、4.3）。


\end{infobox}




\needspace{8\baselineskip}
\section{严格化：下降引理}




\needspace{5\baselineskip}
\subsection{我们到底需要什么假设}



第 2 节用 $o(h)$ 说明了“足够小的步长会下降”，但“足够小”是多大？在多维里怎么判断？要给出可用的结论，需要对函数的光滑程度加一个量化假设。



\textbf{定义（$\beta$-光滑）}：若 $L$ 可微，且存在常数 $\beta>0$ 使得对任意 $\mathbf{w},\mathbf{v}$ 都有


\[
\|\nabla L(\mathbf{w})-\nabla L(\mathbf{v})\|\le\beta\|\mathbf{w}-\mathbf{v}\|
\]


则称 $L$ 是 $\beta$-光滑的。直观上，它说\textbf{梯度不会剧烈变化}，也就是曲面不会突然折起来；$\beta$ 是“曲率的上界”。对二次函数 $L=\frac12\mathbf{w}^\top\mathbf{A}\mathbf{w}$，$\nabla L=\mathbf{A}\mathbf{w}$，可以验证 $\beta=\lambda_{\max}(\mathbf{A})$，即最大特征值。



$\beta$-光滑有一个极其好用的等价形式（对凸函数它是等价的，一般情形下由积分可得）：


\begin{equation*}
L(\mathbf{w}+\Delta)\ \le\ L(\mathbf{w})+\nabla L(\mathbf{w})^\top\Delta+\frac{\beta}{2}\|\Delta\|^2 \tag{4.1}
\end{equation*}


右边是一个\textbf{二次上界}：真正的函数值，永远不超过“切线 + 一个抛物线”。这个不等式就是泰勒展开加上余项控制，把第 2 节的 $o(h)$ 变成了具体的 $\frac{\beta}{2}\|\Delta\|^2$。




\needspace{5\baselineskip}
\subsection{引理与证明}



\begin{goodbox}
{\bfseries\sffamily\small\color{good} 下降引理}\par\vspace{2pt}
设 $L$ 是 $\beta$-光滑的，取 $\Delta=-\eta\nabla L(\mathbf{w})$，则


\end{goodbox}



\textbf{证明}：把 $\Delta=-\eta\nabla L$ 代入 (4.1)，记 $\mathbf{g}=\nabla L(\mathbf{w})$：


\[
\begin{aligned}
L(\mathbf{w}-\eta\mathbf{g})
&\le L(\mathbf{w})+\mathbf{g}^\top(-\eta\mathbf{g})+\frac{\beta}{2}\|-\eta\mathbf{g}\|^2\\[4pt]
&=L(\mathbf{w})-\eta\|\mathbf{g}\|^2+\frac{\beta\eta^2}{2}\|\mathbf{g}\|^2\\[4pt]
&=L(\mathbf{w})-\eta\left(1-\frac{\beta\eta}{2}\right)\|\mathbf{g}\|^2
\end{aligned}
\]


证毕。$\square$




\needspace{5\baselineskip}
\subsection{步长上界的来源}



看括号里的 $1-\frac{\beta\eta}{2}$。只要它为正，即


\[
\eta<\frac{2}{\beta}
\]


右边就是严格下降，而且下降量与 $\eta\|\mathbf{g}\|^2$ 成正比。这一个不等式解释了前面所有的现象：



\begin{itemize}
\item \textbf{为什么 $\eta$ 太大会发散}：$\eta>2/\beta$ 时系数为负，$L$ 反而增大；且每步放大，误差像等比数列一样爆掉。
\item \textbf{为什么一维的临界值是 $2/L''$}：一维时 $\beta=|L''|$，完全吻合第 2 节算出的 $2/14$。
\item \textbf{为什么特征尺度重要}：$\beta$ 是最大曲率；把所有特征放大 10 倍，$\beta$ 放大 100 倍，允许的步长就缩小 100 倍。这就是不标准化时训练容易崩的直接原因。
\item \textbf{为什么取 $\eta=1/\beta$ 是“标准答案”}：此时系数 $1-\frac12=\frac12$，每步至少下降 $\frac{1}{2\beta}\|\mathbf{g}\|^2$，是保证下降的前提下能取的最大步长。
\end{itemize}



\begin{warnbox}
{\bfseries\sffamily\small\color{warn} 但它只是充分条件}\par\vspace{2pt}
$\eta<2/\beta$ 保证下降，不满足它也未必发散（有些方向可能仍然收敛）。反过来，满足它也\textbf{不保证}每步都下降，因为它只约束了一阶近似，实际问题里还有浮点误差、小批量噪声、非光滑点。第 9 节的感知机就是活生生的例子：它的“梯度”是次梯度，下降引理根本不适用，实测有超过一半的更新让总损失上升。


\end{warnbox}




\needspace{5\baselineskip}
\subsection{从“每步下降”到“收敛”}



下降引理只说“每步降一点”，降一点能不能保证最终到达谷底？还需要防止一种情况：步子越来越小，停在半山腰。对凸函数可以排除这种情况。



把引理在 $\eta=1/\beta$ 下写成


\[
L(\mathbf{w}_{k+1})\le L(\mathbf{w}_k)-\frac{1}{2\beta}\|\nabla L(\mathbf{w}_k)\|^2
\]


移项并对 $k=0,\dots,K-1$ 求和，中间项抵消：


\[
\sum_{k=0}^{K-1}\|\nabla L(\mathbf{w}_k)\|^2\ \le\ 2\beta\big[L(\mathbf{w}_0)-L(\mathbf{w}_K)\big]\ \le\ 2\beta\big[L(\mathbf{w}_0)-L^\star\big]
\]


右边是常数，所以左边这个和有限，梯度的平方平均必然趋于零：


\[
\min_{0\le k<K}\|\nabla L(\mathbf{w}_k)\|^2\ \le\ \frac{2\beta\,[L(\mathbf{w}_0)-L^\star]}{K}\ \longrightarrow\ 0
\]


这就是梯度下降的收敛性：\textbf{梯度的最小值以 $O(1/K)$ 的速度趋于零}。若 $L$ 还是凸的，梯度为零意味着全局最优，于是 $L(\mathbf{w}_K)-L^\star=O(1/K)$。



\begin{infobox}
{\bfseries\sffamily\small\color{accent} 这个结论的边界，比结论本身更重要}\par\vspace{2pt}
\begin{itemize}
\item 它保证的是“梯度趋于零”，\textbf{不是}“到达全局最优”。非凸时梯度为零的点可能是鞍点或局部极小。
\item 它保证的是“某个 $k$ 很好”，\textbf{不是}“最后一个 $k$ 很好”，实际中常需要记录历史最优点。
\item $O(1/K)$ 是对最坏情况的上界，实践中常常更快（尤其强凸时是指数收敛）。
\end{itemize}


\end{infobox}




\needspace{8\baselineskip}
\section[二次函数的完整解]{二次函数的完整解\hfill{\normalfont\sffamily\small\color{accent}选读}}



二次函数是唯一能写出\textbf{精确}收敛公式的情形。它值得单独看，因为它把“条件数”“振荡”“收敛速率”这些模糊的词变成了可计算的数。




\needspace{5\baselineskip}
\subsection{坐标变换：解耦}



考虑


\[
L(\mathbf{w})=\frac12\mathbf{w}^\top\mathbf{A}\mathbf{w}-\mathbf{b}^\top\mathbf{w},\qquad \mathbf{A}=\mathbf{A}^\top\succ0
\]


梯度 $\nabla L=\mathbf{A}\mathbf{w}-\mathbf{b}$，最优解满足 $\mathbf{A}\mathbf{w}^\star=\mathbf{b}$。更新式两边减去 $\mathbf{w}^\star$：


\[
\mathbf{w}_{k+1}-\mathbf{w}^\star
=\mathbf{w}_k-\eta(\mathbf{A}\mathbf{w}_k-\mathbf{b})-\mathbf{w}^\star
=(\mathbf{I}-\eta\mathbf{A})(\mathbf{w}_k-\mathbf{w}^\star)
\]


记误差 $\mathbf{e}_k=\mathbf{w}_k-\mathbf{w}^\star$，则 $\mathbf{e}_{k+1}=(\mathbf{I}-\eta\mathbf{A})\mathbf{e}_k$，于是


\[
\mathbf{e}_K=(\mathbf{I}-\eta\mathbf{A})^K\mathbf{e}_0
\]


矩阵的幂不好直接看，但 $\mathbf{A}$ 对称正定，可以正交对角化 $\mathbf{A}=\mathbf{Q}\bm{\Lambda}\mathbf{Q}^\top$，其中 $\mathbf{Q}$ 的列是单位正交特征向量，$\bm{\Lambda}=\mathrm{diag}(\lambda_1,\dots,\lambda_d)$。令 $\mathbf{z}=\mathbf{Q}^\top\mathbf{e}$（相当于把坐标旋转到特征向量方向），则


\[
\mathbf{z}_{k+1}=(\mathbf{I}-\eta\bm{\Lambda})\mathbf{z}_k
\quad\Longrightarrow\quad
z_k^{(i)}=(1-\eta\lambda_i)^k\,z_0^{(i)}
\]


\textbf{在一个“好”的坐标系里，$d$ 维问题拆成了 $d$ 个互不干扰的一维问题}，每个方向独立地按自己的比率 $(1-\eta\lambda_i)$ 衰减。这是整篇文章最漂亮的一步。




\needspace{5\baselineskip}
\subsection{三个结论}



\textbf{（1）收敛条件。} 所有方向都要衰减，需要 $|1-\eta\lambda_i|<1$ 对所有 $i$ 成立，即


\[
0<\eta<\frac{2}{\lambda_{\max}}
\]


与第 4 节的 $\eta<2/\beta$ 一致，因为 $\beta=\lambda_{\max}$。



\textbf{（2）收敛速率由条件数决定。} 最慢的方向是 $\lambda$ 最小（或最大）的那个，其衰减率是 $\max_i|1-\eta\lambda_i|$。可以证明最优步长 $\eta^\star=\frac{2}{\lambda_{\max}+\lambda_{\min}}$，此时最慢衰减率为


\[
\frac{\kappa-1}{\kappa+1},\qquad \kappa=\frac{\lambda_{\max}}{\lambda_{\min}}
\]


$\kappa$ 叫条件数。$\kappa=1$（各方向曲率相同）时衰减率为 0，一步到位；$\kappa=100$ 时衰减率约 $0.98$，要几百步才降一个数量级。\textbf{这就是“狭长山谷里梯度下降很慢”的定量版本}。



\textbf{（3）振荡的根源。} 若某个方向的 $\eta\lambda_i>1$，则 $1-\eta\lambda_i<0$，该方向每步变号，表现为来回横跳。曲率大的方向最容易振荡，因为 $\lambda_i$ 大。实际训练中看到的“之字形”，就是不同方向衰减率差异过大的直接结果。



\begin{examplebox}
{\bfseries\sffamily\small\color{black} 算一遍}\par\vspace{2pt}
取 $L=\frac12(w_1^2+10w_2^2)$，从 $(10,1)$ 出发，$\lambda=(1,10)$，$\beta=10$，稳定条件 $\eta<0.2$。


\begin{itemize}
\item $\eta=0.1$：两方向衰减率 $0.9$ 与 $0$。$w_1$ 每步乘 0.9（慢），$w_2$ 一步到位（快）。所以轨迹几乎是一条沿 $w_1$ 轴缓慢滑行的直线。
\item $\eta=0.15$：衰减率 $0.85$ 与 $-0.5$。$w_2$ 仍然快，但每步变号，轨迹呈小幅锯齿。
\item $\eta=0.19$：衰减率 $0.81$ 与 $-0.9$。$w_2$ 方向几乎不衰减地来回跳，收敛极慢，这就是“贴着稳定边界走”的代价。
\end{itemize}


三个步长都满足 $\eta<0.2$，都没有发散，但收敛速度差出一个数量级。\textbf{“不发散”和“收敛得快”是两件独立的事。}


\end{examplebox}




\needspace{8\baselineskip}
\section{三种梯度下降}



到这里，公式 $\mathbf{w}\leftarrow\mathbf{w}-\eta\nabla L(\mathbf{w})$ 里的梯度，指的是\textbf{在整个训练集上算出的}梯度。但真实数据可能有几百万条，每算一次梯度都要把所有样本过一遍，代价极高。于是有了三种变体。




\needspace{5\baselineskip}
\subsection{批量梯度下降（BGD）}


\[
\nabla L(\mathbf{w})=\frac{1}{N}\sum_{i=1}^{N}\nabla\ell\big(\mathbf{w};\mathbf{x}_i,y_i\big)
\]


每次用全部 $N$ 个样本算梯度。优点是方向准、下降稳定；缺点是单步代价 $O(N)$，数据大时一步都走不动。第 2–5 节分析的都是它。




\needspace{5\baselineskip}
\subsection{随机梯度下降（SGD）}


\[
\mathbf{w}\leftarrow\mathbf{w}-\eta\,\nabla\ell\big(\mathbf{w};\mathbf{x}_{i_k},y_{i_k}\big)
\]


每次只取\textbf{一个}样本 $i_k$（通常按随机顺序轮换），用它的梯度近似整体梯度。单步代价 $O(1)$，同样时间内能更新 $N$ 次。



关键在于：如果样本是均匀随机抽取的，那么单样本梯度是整体梯度的\textbf{无偏估计}：


\[
\mathbb{E}_{i}\big[\nabla\ell(\mathbf{w};\mathbf{x}_i,y_i)\big]=\frac{1}{N}\sum_{i=1}^{N}\nabla\ell(\mathbf{w};\mathbf{x}_i,y_i)=\nabla L(\mathbf{w})
\]


所以“平均而言”它指向正确的方向，但每一步都有噪声，方差约为 $\sigma^2$。这带来两个后果：



\begin{itemize}
\item \textbf{好的一面}：噪声可以帮助跳出较浅的局部极小和鞍点，也带来正则化效果；
\item \textbf{坏的一面}：即使到了最优点附近，梯度也不为零，参数会在最优点周围\textbf{不停抖动}，无法精确停住。因此 SGD 必须配合\textbf{递减的学习率}（如 $\eta_k\propto1/\sqrt{k}$ 或 $1/k$），才能保证收敛。
\end{itemize}




\needspace{5\baselineskip}
\subsection{小批量（mini-batch）}


\[
\mathbf{w}\leftarrow\mathbf{w}-\eta\cdot\frac{1}{B}\sum_{i\in\mathcal{B}_k}\nabla\ell(\mathbf{w};\mathbf{x}_i,y_i)
\]


每次取 $B$ 个样本（如 32、128、256）。这是实践中的默认选择，原因是它同时利用了两件事：



\begin{itemize}
\item \textbf{矩阵运算的效率}：$B$ 个样本的梯度可以用一次矩阵乘法算完，GPU 上 $B=128$ 的耗时和 $B=1$ 相差无几；
\item \textbf{方差与代价的折中}：$B$ 个独立样本平均后，梯度方差降到 $\sigma^2/B$，方向更稳。
\end{itemize}



\begin{center}
\footnotesize
\setlength{\tabcolsep}{3pt}
\begin{tabularx}{\linewidth}{@{}>{\hsize=0.833\hsize\raggedright\arraybackslash}X>{\hsize=0.972\hsize\raggedright\arraybackslash}X>{\hsize=0.694\hsize\raggedright\arraybackslash}X>{\hsize=0.694\hsize\raggedright\arraybackslash}X>{\hsize=1.807\hsize\raggedright\arraybackslash}X@{}}
\toprule
\textbf{方法} & \textbf{单步用的样本} & \textbf{梯度方差} & \textbf{单步代价} & \textbf{典型用途} \\
\midrule
BGD & 全部 $N$ & $0$ & $O(Nd)$ & 小数据、凸问题、理论分析 \\
SGD & 1 & $\sigma^2$ & $O(d)$ & 在线学习、超大样本流 \\
mini-batch & $B$ & $\sigma^2/B$ & $O(Bd)$ & 深度学习默认 \\
\bottomrule
\end{tabularx}
\end{center}



\begin{infobox}
{\bfseries\sffamily\small\color{accent} 收敛速度的差别}\par\vspace{2pt}
凸且光滑时，BGD 的损失以 $O(1/K)$ 下降（$K$ 是\textbf{迭代次数}）；SGD 由于噪声，只能达到 $O(1/\sqrt{K})$。但一次迭代的代价相差 $N$ 倍，所以按\textbf{总计算量}比较时，SGD 往往更快。这是一个“迭代次数”与“单次代价”的经典权衡，不能只看其中一个。


\end{infobox}




\needspace{8\baselineskip}
\section{手算全过程}



理论说完了，现在把它跑一遍。这个例子小到可以在纸上算完，但已经包含了多维梯度下降的所有特征。



目标函数


\[
L(w_1,w_2)=\frac12\left(w_1^2+3w_2^2\right)
\]


梯度


\[
\nabla L=\begin{pmatrix}w_1\\[3pt]3w_2\end{pmatrix}
\]


初始点 $(w_1,w_2)=(2,\,1)$，学习率 $\eta=0.5$。



先检查稳定性：$\mathbf{A}=\mathrm{diag}(1,3)$，$\lambda_{\max}=3$，要求 $\eta<2/3\approx0.667$，$0.5$ 满足。而且第 5 节给出最优步长 $\eta^\star=2/(\lambda_{\max}+\lambda_{\min})=2/4=0.5$，\textbf{正好取到了最优值}。



每一步的更新


\[
\begin{pmatrix}w_1\\w_2\end{pmatrix}_{k+1}
=\begin{pmatrix}w_1\\w_2\end{pmatrix}_{k}
-0.5\begin{pmatrix}w_1\\3w_2\end{pmatrix}_{k}
=\begin{pmatrix}0.5\,w_1\\-0.5\,w_2\end{pmatrix}_{k}
\]


\begin{center}
\footnotesize
\setlength{\tabcolsep}{3pt}
\begin{tabularx}{\linewidth}{@{}>{\hsize=0.889\hsize\raggedleft\arraybackslash}X>{\hsize=0.889\hsize\raggedleft\arraybackslash}X>{\hsize=1.111\hsize\raggedleft\arraybackslash}X>{\hsize=0.889\hsize\raggedleft\arraybackslash}X>{\hsize=1.111\hsize\raggedleft\arraybackslash}X>{\hsize=1.111\hsize\raggedleft\arraybackslash}X@{}}
\toprule
\textbf{$k$} & \textbf{$w_1$} & \textbf{$w_2$} & \textbf{$\partial L/\partial w_1=w_1$} & \textbf{$\partial L/\partial w_2=3w_2$} & \textbf{$L$} \\
\midrule
0 & 2 & 1 & 2 & 3 & 3.5 \\
1 & 1 & −0.5 & 1 & −1.5 & 0.875 \\
2 & 0.5 & 0.25 & 0.5 & 0.75 & 0.21875 \\
3 & 0.25 & −0.125 & 0.25 & −0.375 & 0.054688 \\
4 & 0.125 & 0.0625 & 0.125 & 0.1875 & 0.013672 \\
5 & 0.0625 & −0.03125 & 0.0625 & −0.09375 & 0.003418 \\
\bottomrule
\end{tabularx}
\end{center}
{\footnotesize\color{gray}\raggedright 最优解是 $(0,0)$，最优值 $L^\star=0$。第 0 行是初值，梯度在更新\textbf{之前}计算。\par}



读这张表，能看出四件事：



\begin{enumerate}
\item \textbf{$w_1$ 严格单调减半}：$2\to1\to0.5\to0.25\cdots$，因为 $1-\eta\lambda_1=1-0.5=0.5$。
\item \textbf{$w_2$ 每步变号并且减半}：$1\to-0.5\to0.25\to-0.125\cdots$，因为 $1-\eta\lambda_2=1-1.5=-0.5$。负号就是“跨过谷底”的直接证据。
\item \textbf{$L$ 每步恰好缩小到四分之一}：因为两个方向都乘了 $0.5$，二次型按 $0.5^2=0.25$ 缩放。这正好验证第 5 节的衰减率 $(\kappa-1)/(\kappa+1)=2/4=0.5$，以及 $\kappa=3$。
\item \textbf{梯度的大小随 $L$ 同步缩小}：越接近谷底，步子自动越小。这是梯度下降“自动减速”的机制，也是它不需要手动调整步长衰减的原因（凸二次情形）。
\end{enumerate}



\begin{infobox}
{\bfseries\sffamily\small\color{accent} 注意第 2 点}\par\vspace{2pt}
$w_2$ 在 $1$ 和 $-0.5$ 之间来回跳，看起来像是“错了”，其实完全正确：$w_2$ 的曲率是 3，最优步长的两倍正好等于它的周期，所以每步都精确落在对面。这种“振荡但收敛”是梯度下降的正常行为，不要误判为发散。判断标准只有一个：\textbf{误差的绝对值是不是在缩小}。


\end{infobox}




\needspace{8\baselineskip}
\section{收敛性到底保证了什么}



“梯度下降会收敛”这句话，在不同假设下含义完全不同。这一节把三种情形分开说清楚，因为混用它们是最常见的理论误读。




\needspace{5\baselineskip}
\subsection{凸且光滑：$O(1/K)$}



若 $L$ 凸、$\beta$-光滑，取 $\eta=1/\beta$，则


\[
L(\mathbf{w}_K)-L^\star\ \le\ \frac{\beta\|\mathbf{w}_0-\mathbf{w}^\star\|^2}{2K}=O\!\left(\frac{1}{K}\right)
\]


含义：迭代次数增加 10 倍，误差降到约 1/10。想再降一个数量级，就要再花 10 倍步数。\textbf{这是“次线性”收敛}，慢，但对一般凸问题已经是最优的一阶方法速率（有下界证明）。




\needspace{5\baselineskip}
\subsection{强凸且光滑：线性收敛}



若还满足 $\mu$-强凸（曲率有下界 $\mu>0$），则


\[
L(\mathbf{w}_K)-L^\star\ \le\ \left(1-\frac{\mu}{\beta}\right)^{K}\big(L(\mathbf{w}_0)-L^\star\big)
\]


误差按固定比例缩小，叫\textbf{线性收敛}（在数值分析里“线性”指误差等比，不是指函数是线性的）。速率由条件数 $\kappa=\beta/\mu$ 决定：$\kappa$ 越大越慢。第 5 节的二次函数正是这种情形，那时 $\mu=\lambda_{\min}$、$\beta=\lambda_{\max}$，与 $(\kappa-1)/(\kappa+1)$ 的结论一致（量级相同）。




\needspace{5\baselineskip}
\subsection{非凸：只能到驻点}



神经网络的损失几乎都是非凸的。此时能保证的只有 4.4 节那条：


\[
\min_{k<K}\|\nabla L(\mathbf{w}_k)\|^2\ \le\ O\!\left(\frac{1}{K}\right)
\]


即\textbf{存在某一步梯度很小}。而“梯度小”可能是极小点，也可能是鞍点，还可能是平坦的高原。理论上无法区分，实践中靠随机初始化、噪声、动量等手段逃离。



\begin{warnbox}
{\bfseries\sffamily\small\color{warn} "收敛"不等于"找到最优"}\par\vspace{2pt}
非凸时，梯度下降甚至可能收敛到一个\textbf{比初始点还差的局部极小}。这不是实现 bug，而是问题本身的性质。深度学习中它之所以不致命，是因为高维空间的局部极小往往接近全局最优，真正麻烦的是鞍点和大片平坦区域。


\end{warnbox}




\needspace{5\baselineskip}
\subsection{一个特殊情形：感知机}



你在学感知机，它正好是一个“看起来不该收敛、却能被证明收敛”的例子。它的损失


\[
L(\mathbf{w},b)=-\sum_{i\in M}y_i\left(\mathbf{w}^\top\mathbf{x}_i+b\right)
\]


是\textbf{分段线性}的（不是光滑的），而且 $M$ 随参数变化，所以第 4 节的下降引理不适用，实测中确实会出现“某次更新让总损失上升”。



但它有另一个结论，来自 Novikoff 定理：若数据线性可分，且 $\|\hat{\mathbf{x}}_i\|\le R$、存在单位向量 $\mathbf{u}$ 使 $y_i\mathbf{u}^\top\hat{\mathbf{x}}_i\ge\gamma$，则误分类次数


\[
k\ \le\ \left(\frac{R}{\gamma}\right)^2
\]


这里 $\hat{\mathbf{x}}=(\mathbf{x},1)$ 是把偏置吸收进去的增广向量。\textbf{它保证的是“有限步停止”，不是“损失单调下降”}。两种收敛性（函数值下降 vs 迭代终止）在非光滑问题里会分家，这是很好的提醒：不要用一个定理的结论去套另一个问题的假设。




\needspace{8\baselineskip}
\section{三个真实模型的推导}



梯度下降的公式是通用的，但每个模型的梯度长得不一样。这一节把三个最基础的模型完整推一遍，你会发现它们最后都落到同一个更新形式上。




\needspace{5\baselineskip}
\subsection{线性回归（平方损失）}



\textbf{模型}：$\hat{y}=\mathbf{w}^\top\mathbf{x}$（为简洁省略截距，把它并进 $\mathbf{x}$ 即可）。



\textbf{损失}（均方误差的一半，$N$ 个样本）：


\[
L(\mathbf{w})=\frac{1}{2N}\sum_{i=1}^{N}\left(\mathbf{w}^\top\mathbf{x}_i-y_i\right)^2
\]


\textbf{求梯度}：记残差 $r_i=\mathbf{w}^\top\mathbf{x}_i-y_i$。对单个样本，用链式法则：


\[
\frac{\partial}{\partial\mathbf{w}}\frac12 r_i^2=r_i\cdot\frac{\partial r_i}{\partial\mathbf{w}}=r_i\,\mathbf{x}_i
\]


（这里用到 $\partial(\mathbf{w}^\top\mathbf{x}_i)/\partial\mathbf{w}=\mathbf{x}_i$，可以逐分量验证：$\partial(\sum_j w_jx_{ij})/\partial w_k=x_{ik}$。）对所有样本求和：


\[
\nabla L(\mathbf{w})=\frac{1}{N}\sum_{i=1}^{N}r_i\mathbf{x}_i=\frac{1}{N}\mathbf{X}^\top(\mathbf{X}\mathbf{w}-\mathbf{y})
\]


其中 $\mathbf{X}\in\mathbb{R}^{N\times d}$ 每行是一个样本。\textbf{更新}：


\[
\mathbf{w}\leftarrow\mathbf{w}-\eta\cdot\frac{1}{N}\mathbf{X}^\top(\mathbf{X}\mathbf{w}-\mathbf{y})
\]


\begin{infobox}
{\bfseries\sffamily\small\color{accent} 一个容易忽略的点}\par\vspace{2pt}
对线性回归，令梯度为零可以直接解出 $\mathbf{w}^\star=(\mathbf{X}^\top\mathbf{X})^{-1}\mathbf{X}^\top\mathbf{y}$，这就是正规方程。既然有解析解，为什么还用梯度下降？因为 $d$ 很大时求逆是 $O(d^3)$、存 $\mathbf{X}^\top\mathbf{X}$ 要 $O(d^2)$ 内存，而梯度下降每步只要 $O(Nd)$、内存 $O(d)$。\textbf{算法选择取决于规模，不是取决于哪个“更高级”。}


\end{infobox}




\needspace{5\baselineskip}
\subsection{逻辑回归（交叉熵）}



\textbf{模型}：$p=\sigma(z)$，其中 $z=\mathbf{w}^\top\mathbf{x}$，$\sigma(z)=\dfrac{1}{1+e^{-z}}$ 是 sigmoid。



\textbf{损失}（负对数似然，标签 $y\in\{0,1\}$）：


\[
L(\mathbf{w})=-\frac{1}{N}\sum_{i=1}^{N}\Big[y_i\log\sigma(z_i)+(1-y_i)\log\big(1-\sigma(z_i)\big)\Big]
\]


\textbf{关键性质}：sigmoid 的导数有一个漂亮的形式，先算出来：


\[
\sigma'(z)=\frac{e^{-z}}{(1+e^{-z})^2}=\sigma(z)\big(1-\sigma(z)\big)
\]


对单个样本求导，记 $p_i=\sigma(z_i)$：


\[
\frac{\partial \ell_i}{\partial z_i}=-\frac{y_i}{p_i}\cdot p_i(1-p_i)+\frac{1-y_i}{1-p_i}\cdot p_i(1-p_i)=-\,y_i(1-p_i)+(1-y_i)p_i=p_i-y_i
\]


再乘上 $\partial z_i/\partial\mathbf{w}=\mathbf{x}_i$：


\[
\nabla L(\mathbf{w})=\frac{1}{N}\sum_{i=1}^{N}(p_i-y_i)\,\mathbf{x}_i=\frac{1}{N}\mathbf{X}^\top(\mathbf{p}-\mathbf{y})
\]


注意它的形状与线性回归\textbf{完全一致}：都是“残差 $\times$ 特征”再平均。区别只在残差的定义：线性回归用 $\hat{y}-y$，逻辑回归用 $p-y$。这不是巧合，而是广义线性模型的结构。




\needspace{5\baselineskip}
\subsection{感知机（合页型损失）}



\textbf{模型}：$f(\mathbf{x})=\operatorname{sign}(\mathbf{w}^\top\mathbf{x}+b)$，标签 $y\in\{-1,+1\}$。



\textbf{损失}：只统计被分错的点，


\[
L(\mathbf{w},b)=-\sum_{i\in M}y_i\left(\mathbf{w}^\top\mathbf{x}_i+b\right),\qquad M=\{i:\ y_i(\mathbf{w}^\top\mathbf{x}_i+b)\le0\}
\]


\textbf{求梯度}：对 $M$ 内每个点，$y_i(\mathbf{w}^\top\mathbf{x}_i+b)$ 对 $\mathbf{w}$ 的导数是 $y_i\mathbf{x}_i$，前面有负号：


\[
\nabla_{\mathbf{w}}L=-\sum_{i\in M}y_i\mathbf{x}_i,\qquad \frac{\partial L}{\partial b}=-\sum_{i\in M}y_i
\]


\textbf{更新}：减梯度等于加


\[
\mathbf{w}\leftarrow\mathbf{w}+\eta\sum_{i\in M}y_i\mathbf{x}_i,\qquad b\leftarrow b+\eta\sum_{i\in M}y_i
\]


教材里常见的写法是每次只取一个误分类点（随机梯度）：


\[
\mathbf{w}\leftarrow\mathbf{w}+\eta\,y_i\mathbf{x}_i,\qquad b\leftarrow b+\eta\,y_i
\]


\begin{warnbox}
{\bfseries\sffamily\small\color{warn} 为什么这里不能套用第 4 节的定理}\par\vspace{2pt}
$L$ 在 $y_i(\mathbf{w}^\top\mathbf{x}_i+b)=0$ 处有折点，不可微；梯度只在严格误分类（$y_if_i<0$）时存在。此时更新方向是\textbf{次梯度}，下降引理不成立，“每步都下降”是错的。你在上一轮看到的统计——28 次更新中有 15 次总损失上升——正是这个原因。


\end{warnbox}




\needspace{5\baselineskip}
\subsection{三者对照}



\begin{center}
\footnotesize
\setlength{\tabcolsep}{3pt}
\begin{tabularx}{\linewidth}{@{}>{\hsize=1.053\hsize\raggedright\arraybackslash}X>{\hsize=0.842\hsize\raggedright\arraybackslash}X>{\hsize=0.842\hsize\raggedright\arraybackslash}X>{\hsize=1.263\hsize\raggedright\arraybackslash}X@{}}
\toprule
\textbf{模型} & \textbf{损失} & \textbf{梯度} & \textbf{更新} \\
\midrule
线性回归 & $\frac{1}{2N}\|\mathbf{X}\mathbf{w}-\mathbf{y}\|^2$ & $\frac1N\mathbf{X}^\top(\mathbf{X}\mathbf{w}-\mathbf{y})$ & $\mathbf{w}-\eta\cdot$ 梯度 \\
逻辑回归 & 交叉熵 & $\frac1N\mathbf{X}^\top(\mathbf{p}-\mathbf{y})$ & $\mathbf{w}-\eta\cdot$ 梯度 \\
感知机 & $-\sum_{M}y_i f_i$ & $-\sum_{M}y_i\mathbf{x}_i$ & $\mathbf{w}+\eta\sum_{M}y_i\mathbf{x}_i$ \\
\bottomrule
\end{tabularx}
\end{center}
{\footnotesize\color{gray}\raggedright 前两行是“残差 \ensuremath{\times} 特征”，第三行因为损失自带负号，更新变成加法。三者都是同一个公式的特例。\par}




\needspace{8\baselineskip}
\section{八个常见误区}



这些说法在教材和博客里反复出现，但每一句都需要加上限定条件才成立。




\needspace{5\baselineskip}
\subsection*{误区 1：梯度下降每一步都会让损失变小}



只在可微、$\eta<2/\beta$、且用真实梯度的条件下成立。小批量有噪声、感知机不可微、步长过大时都可能上升。\textbf{正确说法}：在满足条件的凸光滑问题上，损失单调不增。




\needspace{5\baselineskip}
\subsection*{误区 2：负梯度是“下降最快的方向”}



只在\textbf{无穷小步长 + 欧氏长度}下成立。第 3.3 节的数值表里，夹角 89° 的方向一阶预测还有微小下降，实测却上升了。有限步长时，二阶项可以改变排序。




\needspace{5\baselineskip}
\subsection*{误区 3：梯度为零就是最小值}



可能是极大值、鞍点，或者平坦区域。一维反例 $f(w)=w^3$ 在 $w=0$ 处梯度为零，但那既不是极大也不是极小。高维中鞍点远比局部极小常见。




\needspace{5\baselineskip}
\subsection*{误区 4：学习率越小越好}



太小会导致收敛极慢、在平坦区停滞，还可能因浮点精度而“假收敛”（更新量小于机器精度）。正确做法是选在稳定上界内、配合衰减策略。




\needspace{5\baselineskip}
\subsection*{误区 5：损失下降就说明模型变好了}



训练损失下降只说明\textbf{拟合训练集}更好，可能是过拟合。要同时看验证集。梯度下降优化的是训练损失，不直接优化泛化能力。




\needspace{5\baselineskip}
\subsection*{误区 6：特征不用标准化}



标准化会改变 $\beta$ 和 $\mu$，从而改变稳定步长和条件数。某列量纲是万元、另一列是比例时，梯度会被大尺度特征主导，收敛极慢甚至发散。\textbf{标准化不是可选的美化步骤，而是算法能正常工作的前提之一。}




\needspace{5\baselineskip}
\subsection*{误区 7：SGD 只是“更快的梯度下降”}



它更快是因为单步便宜，但收敛速率更差（$O(1/\sqrt{K})$ 对 $O(1/K)$），而且不会精确停在最优点。它的价值在于\textbf{单位计算量下的进展}，以及噪声带来的逃离鞍点能力。




\needspace{5\baselineskip}
\subsection*{误区 8：感知机的学习率会影响收敛速度}



不会。$\eta$ 只把 $(\mathbf{w},b)$ 整体缩放，分类判定只看符号，所以更新序列与 $\eta$ 无关，迭代次数完全相同（实测 $\eta=0.01,1,100$ 都是 8 次）。这是感知机损失\textbf{正齐次}的特殊后果，不能推广到别的模型。




\needspace{8\baselineskip}
\section{自测}



合上讲义，先自己算，再展开对照。




\needspace{5\baselineskip}
\subsection*{题 1}



$f(w)=w^3$，从 $w_0=1$ 出发，$\eta=0.1$，沿负梯度走一步后 $w_1$ 是多少？$f$ 下降了吗？如果继续走，会停在哪里？



\begin{answerbox}
{\bfseries\sffamily\small\color{accent} 解答}\par\vspace{2pt}
$f'(w)=3w^2$，$f'(1)=3$，所以 $w_1=1-0.1\times3=0.7$。$f(1)=1$，$f(0.7)=0.343$，确实下降。


但 $w=0$ 处 $f'=0$，梯度下降会停在那里——而它不是极小值（左边更小）。这说明“梯度为零”只是\textbf{驻点}，不保证最优。这个例子还说明：非凸问题上，梯度下降的结果强烈依赖初始点。


\end{answerbox}




\needspace{5\baselineskip}
\subsection*{题 2}



$L(w)=2w^2$，从 $w_0=1$ 出发。分别取 $\eta=0.1$、$0.2$、$0.3$、$0.6$，写出 $w_k$ 的通项，判断是否收敛。



\begin{answerbox}
{\bfseries\sffamily\small\color{accent} 解答}\par\vspace{2pt}
这里 $\lambda=L''=4$，$w^\star=0$，误差递推 $e_{k+1}=(1-4\eta)e_k$。


\begin{itemize}
\item $\eta=0.1$：$1-0.4=0.6$，$w_k=0.6^k$，单调收敛。
\item $\eta=0.2$：$1-0.8=0.2$，$w_k=0.2^k$，收敛更快。
\item $\eta=0.3$：$1-1.2=-0.2$，$w_k=(-0.2)^k$，收敛但每步变号。
\item $\eta=0.6$：$1-2.4=-1.4$，$|{-1.4}|>1$，发散。
\end{itemize}


临界值 $2/\lambda=0.5$。注意 $\eta=0.3$ 仍收敛，只是振荡——\textbf{振荡不等于发散}，判断依据是 $|1-\eta\lambda|$ 是否小于 1。


\end{answerbox}




\needspace{5\baselineskip}
\subsection*{题 3}



$L=\frac12(w_1^2+100w_2^2)$，用梯度下降会有什么现象？如何改进（至少说两种）？



\begin{answerbox}
{\bfseries\sffamily\small\color{accent} 解答}\par\vspace{2pt}
条件数 $\kappa=100$。为保证 $w_2$ 方向不振荡，需要 $\eta<2/100=0.02$；但此时 $w_1$ 方向的衰减率是 $1-0.02=0.98$，每步只缩小 2\%，要几百步才能收敛。这就是“一个方向被另一个方向绑住”的典型现象。


改进：① 标准化特征，让两列尺度接近，$\kappa$ 变小；② 用牛顿法或拟牛顿法，方向改为 $-\mathbf{H}^{-1}\nabla L$；③ 用动量、Adam 等自适应方法，按方向调整步长；④ 坐标缩放/预条件。任答两种即可。


\end{answerbox}




\needspace{5\baselineskip}
\subsection*{题 4}



感知机的损失在 $y_i(\mathbf{w}^\top\mathbf{x}_i+b)=0$ 处可微吗？这对“每步下降”的结论有什么影响？



\begin{answerbox}
{\bfseries\sffamily\small\color{accent} 解答}\par\vspace{2pt}
不可微。该点左侧（函数间隔为负）斜率是 $-y_i\mathbf{x}_i$，右侧（间隔为正，该点已不在 $M$ 中）斜率是 $0$，左右导数不相等，是一个折点。


因此只能用\textbf{次梯度}（两侧斜率的凸组合）。次梯度法不保证每步下降，只保证在合适步长下“到最优解的距离不增”之类的弱结论；实测中确实会出现总损失上升。


\end{answerbox}




\needspace{5\baselineskip}
\subsection*{题 5}



用自己的话解释：为什么“沿负梯度走”能让所有参数同时下降，而别的方向不行？



\begin{answerbox}
{\bfseries\sffamily\small\color{accent} 解答}\par\vspace{2pt}
一阶变化可以按参数拆成 $\sum_j(\partial L/\partial w_j)\Delta w_j$。只有取 $\Delta w_j=-\eta\,\partial L/\partial w_j$ 时，每一项都等于 $-\eta(\partial L/\partial w_j)^2\le0$，没有一项为正。


换任何其他方向，总有些 $\Delta w_j$ 与对应偏导同号，那些参数就在增加损失，需要靠别的参数补偿，净下降必然更少。这就是“无内耗”的含义，也是 Cauchy–Schwarz 不等式等号条件的直观版本。


\end{answerbox}




\clearpage

\needspace{8\baselineskip}
\section{代码附录}



一个不依赖任何机器学习库的实现，把本文所有公式落到实处。它同时支持批量和小批量，并记录每步的损失，方便你自己验证第 7 节的表格。



\noindent\begin{minipage}{\linewidth}
\begin{lstlisting}[basicstyle=\ttfamily\scriptsize,aboveskip=6pt]
import numpy as np

class GradientDescent:
    """最小二乘 / 一般可微损失的梯度下降。"""

    def __init__(self, lr=0.1, epochs=100, batch_size=None, seed=0):
        self.lr = lr
        self.epochs = epochs
        self.batch_size = batch_size      # None 表示批量梯度下降
        self.seed = seed

    def fit(self, X, y):
        X, y = np.asarray(X, float), np.asarray(y, float)
        n, d = X.shape
        w = np.zeros(d)
        self.history_ = []
        rng = np.random.default_rng(self.seed)

        for _ in range(self.epochs):
            idx = rng.permutation(n) if self.batch_size else np.arange(n)
            bs = self.batch_size or n
            for s in range(0, n, bs):
                b = idx[s:s + bs]
                grad = X[b].T @ (X[b] @ w - y[b]) / len(b)   # 第 9.1 节
                w -= self.lr * grad
                self.history_.append(0.5 * np.mean((X @ w - y) ** 2))
        self.w_ = w
        return self

    def predict(self, X):
        return np.asarray(X, float) @ self.w_


if __name__ == "__main__":
    # 复现第 7 节的二次函数：L = 0.5*(w1^2 + 3*w2^2)
    A = np.diag([1.0, 3.0])
    w = np.array([2.0, 1.0])
    eta = 0.5
    for k in range(6):
        L = 0.5 * w @ A @ w
        print(f"k={k}  w=({w[0]:.4f}, {w[1]:.4f})  L={L:.6f}")
        w = w - eta * (A @ w)
\end{lstlisting}
\end{minipage}



运行输出应当与第 7 节的表格逐行吻合：



\noindent\begin{minipage}{\linewidth}
\begin{lstlisting}[basicstyle=\ttfamily\scriptsize,aboveskip=6pt]
k=0  w=(2.0000, 1.0000)  L=3.500000
k=1  w=(1.0000, -0.5000)  L=0.875000
k=2  w=(0.5000, 0.2500)  L=0.218750
k=3  w=(0.2500, -0.1250)  L=0.054688
k=4  w=(0.1250, 0.0625)  L=0.013672
k=5  w=(0.0625, -0.0313)  L=0.003418
\end{lstlisting}
\end{minipage}



\begin{goodbox}
{\bfseries\sffamily\small\color{good} 怎么用它检验自己的理解}\par\vspace{2pt}
改三个地方试：把 \texttt{eta} 改成 0.7（超过 $2/3$）看何时发散；把 \texttt{A} 改成 \texttt{diag([1, 100])} 看两个方向的收敛速度差异；把 $L$ 改成 $w^3$ 看它停在 $w=0$。这三个实验对应第 5、8、10 节的全部结论。


\end{goodbox}




\par\medskip\noindent\rule{\linewidth}{0.4pt}\medskip\par




\needspace{5\baselineskip}
\subsection*{结语：这份推导的骨架}



把全文压缩成五行，就是梯度下降的全部逻辑：



\begin{enumerate}
\item \textbf{目标}：找一个 $\mathbf{w}$ 使 $L(\mathbf{w})$ 最小。除了极少数情形，没有解析解。
\item \textbf{信息}：只有当前点的梯度 $\nabla L$，它给出每个参数的一阶敏感度。
\item \textbf{方向}：$-\nabla L$ 是唯一让所有参数同向下坡的方向（Cauchy–Schwarz 等号条件）。
\item \textbf{步长}：必须满足 $\eta<2/\beta$，否则一阶近似失效、误差被放大。
\item \textbf{保证}：凸时损失以 $O(1/K)$ 下降；非凸时只能到驻点；不可微时连“每步下降”都没有。
\end{enumerate}



再往深处走，有三个自然的岔路：\textbf{二阶方法}（用 Hessian 修正方向，解决条件数）、\textbf{自适应步长}（Adam、RMSProp，解决各方向尺度不一）、\textbf{随机优化理论}（噪声、方差缩减、非凸景观）。它们的出发点都是这份推导里某一条假设不够用。



如果你想把这条线接到感知机上，最值得自己推一遍的是\textbf{对偶形式}：把 $\mathbf{w}=\sum_i\alpha_iy_i\mathbf{x}_i$ 代回，更新从“改 $\mathbf{w}$”变成“给 $\alpha_i$ 加 $\eta$”，整个过程只依赖 Gram 矩阵 $\mathbf{x}_i^\top\mathbf{x}_j$。这一步会解释为什么感知机能自然核化，也会让你看清“参数空间”和“样本空间”的对应关系。



\end{document}
